\def\title{Is the Hot Hand Real? Evidence from a Permutation and Hierarchical-Based Analysis}
\def\author{Cameron An}
\def\degreemonth{June}
\def\degreeyear{2026}
\def\degree{Master of Science in Statistics}
\def\numberofmembers{3}
   \def\chair{Kevin Ross, Ph.D. \linebreak Professor of Statistics}
   \def\othermemberA{Bret Holladay, Ph.D. \linebreak Assistant Professor of Statistics}
   \def\othermemberB{Immanuel Williams, Ph.D. \linebreak Assistant Professor of Statistics}
\def\keywords{Hot hand, Streak shooting, NBA, Permutation test, Logistic hierarchical model.} % Make sure to end this in a . (i.e. Some, Thing, Interesting.)

\def\abstract{%
The hot-hand phenomenon, often described as the tendency for individuals to experience prolonged streaks of success that exceed what 
would be expected under random performance, has been widely studied across many disciplines, particularly basketball. Early studies 
attempted to evaluate this effect through various techniques, often concluding that the hot-hand phenomenon was largely a myth. However, 
recent studies have begun to revisit previous analyses using improved statistical techniques, with some claiming evidence of a discernible 
hot-hand effect. This study examines the presence of the hot-hand effect in the modern NBA by testing whether observed shooting patterns 
deviate from those simulated under independence. A permutation-based analysis, which preserves observed shot outcomes while randomizing 
order, reveals that the observed shot sequences generally do not exhibit elevated streak behavior, providing limited evidence of a 
hot-hand effect beyond what would be expected by chance. To further account for player and season-level heterogeneity, a logistic 
hierarchical model was constructed to control for in-game variables that contribute to shot success. The estimated effects of being 
in a hot state are small and reinforce the findings of the permutation test after adjusting for additional covariates. 
}

\def\tocpagelegnth{2} % Number of pages your Table of Contents is (used to set correct page numbers for List of Tables and List of Figures)
\def\totflag{1} % 1 if you have at least 1 table, otherwise 0 (adds Table of Tables page)
\def\tofflag{1} % 1 if you have at least 1 figure, otherwise 0 (adds Table of Figures page)